\documentclass[12pt]{article}
\usepackage[margin=0.72in]{geometry}
\usepackage{lmodern}
\usepackage{microtype}
\usepackage{multicol}
\usepackage{fancyhdr}
\usepackage{titlesec}
\usepackage{enumitem}
\usepackage{xcolor}
\usepackage[hidelinks]{hyperref}

\setlength{\columnsep}{0.30in}
\setlength{\parindent}{1.05em}
\setlength{\parskip}{0.24em}
\setlength{\headheight}{15pt}
\titleformat{\section}{\large\bfseries\scshape}{}{0pt}{}
\titleformat{\subsection}{\normalsize\bfseries}{}{0pt}{}
\pagestyle{fancy}
\fancyhf{}
\fancyhead[L]{Cicero on Fate}
\fancyhead[R]{Reading Handout}
\fancyfoot[C]{\thepage}
\renewcommand{\headrulewidth}{0.4pt}

\newcommand{\focusbox}[1]{%
\vspace{0.4em}\noindent\colorbox{gray!12}{\begin{minipage}{0.96\linewidth}#1\end{minipage}}\vspace{0.5em}}

\begin{document}
\begin{center}
{\LARGE\bfseries Are We Free, or Are We Ruled by Fate?}\par\vspace{0.25em}
{\large\scshape Reading 2: Cicero on Fate and Responsibility}\par\vspace{0.45em}
\rule{0.82\textwidth}{0.4pt}\par\vspace{0.7em}
\end{center}

\section*{Vocabulary Preview}
\begin{multicols}{2}
\begin{itemize}[leftmargin=*, itemsep=0.18em]
\item \textbf{Fate}: The idea that events are fixed beforehand or bound by an unavoidable order.
\item \textbf{Necessity}: Something that cannot be otherwise.
\item \textbf{Causation}: The relation between a cause and what follows from it.
\item \textbf{Chance}: What happens without being planned or fixed by a necessary order.
\item \textbf{Stoics}: Ancient philosophers who often taught that the world is ruled by divine reason and fate.
\item \textbf{Epicureans}: Ancient philosophers who denied strict fate and emphasized chance, pleasure, and freedom from fear.
\item \textbf{Assent}: The mind's act of accepting an impression or judgment as true.
\item \textbf{Divination}: The attempt to know future events through signs, omens, or prophecy.
\item \textbf{Prediction}: A statement about what will happen in the future.
\item \textbf{Character}: The settled shape of a person's habits, desires, and judgments.
\item \textbf{Academic}: In Cicero's case, a skeptical philosophical method that argues both sides and follows what seems most probable.
\item \textbf{Moral responsibility}: Being rightly subject to praise, blame, reward, or punishment.
\end{itemize}
\end{multicols}

\section*{Introduction}
Marcus Tullius Cicero was a Roman statesman, orator, and philosopher who lived from 106 to 43 B.C. He wrote \textit{On Fate} near the end of the Roman Republic, during a time of civil war, political danger, and uncertainty about Rome's future.

Cicero inherited a Greek debate. Some Stoic philosophers argued that everything happens by fate: each event follows from earlier causes in a fixed chain. Epicureans rejected this by saying that atoms sometimes ``swerve,'' allowing room for chance and freedom. Cicero is not satisfied with either easy answer. He wants to know whether fate destroys responsibility, and whether human beings can still be praised or blamed if causes shape what they do.

\focusbox{\textbf{Source note:} Adapted and modernized from Cicero, \textit{On Fate} (\textit{De Fato}), selected sections. This classroom version condenses the surviving argument, preserves Cicero's major examples and distinctions, and uses modern student-facing English.}

\begin{multicols}{2}
\section*{Reading}

\subsection*{1. Why Fate Is a Serious Question}
The question of fate is not a small puzzle. It touches almost every part of philosophy. It concerns character, because praise and blame depend on whether actions are our own. It concerns logic, because people make statements about future events: ``This will happen,'' or ``That will not happen.'' It concerns nature, because every event seems to have causes.

If every future event is already fixed, then it seems that human effort may be useless. But if nothing is fixed, then the world may seem ruled by accident. Cicero wants a wiser answer than either despair or confusion.

He asks: when something happens, must we say it happened by fate? Or can some things happen by nature, some by chance, and some by human choice?

\subsection*{2. Nature and Chance Are Not the Same as Fate}
Some people point to strange stories and call them fate. A man is warned to beware a horse and later dies in a place named after a horse. A sailor survives shipwreck but later dies by water. A rock falls in a cave and kills the man who took shelter there. These stories seem impressive because the ending fits the warning.

But Cicero asks whether we need fate to explain them. If there were no word \textit{fate}, and no mysterious force called fate, would the events look different? Rocks would still fall. Storms would still happen. People would still make dangerous choices. Accidents would still occur.

So Cicero does not want to call every surprising event fate. Many things can be explained by nature and chance without claiming that they were fixed beforehand by necessity.

\subsection*{3. Causes Influence Us, but Do They Rule Us?}
Cicero admits that human beings are influenced by causes outside themselves. Climate, health, family, custom, education, and natural temperament all matter. A city with clear air may sharpen the mind. A dense climate may produce sturdier bodies. One person may be quick-witted; another may be slow. One may be naturally calm; another may be easily stirred.

But Cicero denies that such influences determine every choice. The air of Athens may help a student think clearly, but it does not force him to choose one teacher rather than another. A person's natural strength may make athletic victory easier, but it does not by itself decide which race he enters or how he trains.

Some traits are given to us. But our actions are not simply the same thing as those traits. A person may be inclined toward anger without being forced to strike. A person may desire pleasure without being forced to betray a friend. Causes press on us, but Cicero doubts that they erase the difference between influence and necessity.

\subsection*{4. The Stoic Chain of Causes}
The Stoics argued that everything happens through an unbroken chain of causes. Nothing occurs without a cause; every cause has causes before it; therefore the future follows from the past. If this is true, then fate is not a magical prediction. It is the whole order of nature, joined together from beginning to end.

Cicero sees the strength of this view. It explains why events are not random. It also fits ordinary reasoning: if something happens, we usually look for what caused it.

But he worries about the conclusion. If every act of the will is fixed by earlier causes, what becomes of deliberation? Why praise courage or blame betrayal? Why give advice, make laws, teach children, or punish criminals? If all actions are necessary, then no one could have acted otherwise.

\subsection*{5. The Lazy Argument}
Cicero considers a dangerous argument sometimes called the lazy argument.

Suppose someone is sick. If fate has fixed that he will recover, he will recover whether he calls the doctor or not. If fate has fixed that he will not recover, he will not recover whether he calls the doctor or not. Therefore, the lazy person says, calling the doctor is useless.

Cicero rejects this. Even if an outcome is connected to fate, the means may be connected too. If a person is fated to recover, perhaps he is fated to recover by calling the doctor, taking medicine, and obeying treatment. The action is not automatically useless just because the outcome has causes.

This matters for moral life. A student cannot say, ``If I am fated to understand, I will understand whether I study or not.'' Study may be part of the path by which understanding comes. Causes do not always cancel effort; often effort is one of the causes.

\subsection*{6. Assent and the Rolling Cylinder}
The Stoic Chrysippus tried to defend responsibility while keeping fate. He used an image.

Imagine pushing a cylinder. The push begins its movement, but the cylinder rolls because of its own round shape. In the same way, outside events strike the mind and give impressions. But the mind's assent comes from its own character.

Cicero takes this image seriously. It shows that not all causes work in the same way. An outside event may present an opportunity, insult, temptation, or danger. Yet the person's response depends on the kind of soul he has. One person hears an insult and laughs; another hears the same insult and becomes violent. The outside cause is similar, but the inner character differs.

Still, Cicero presses the problem: if the person's character is itself part of the chain of causes, then is assent truly free? If the cylinder rolls because it is shaped that way, and its shape was already fixed, have we really saved responsibility?

\subsection*{7. Fate, Prediction, and Human Choice}
Cicero also considers predictions. If someone truly predicts a future event, does that make the event necessary? If it was true yesterday that a person would act unjustly today, could he avoid acting unjustly?

Cicero does not want truth about the future to become a prison. A statement may turn out true because of what a person freely does, not because the statement forced him to do it. Knowing or predicting an action is not the same as causing it.

This will matter later when students read tragedies about prophecy. If Oedipus or Macbeth hears a prediction, the prediction may shape his imagination and choices. But does it force the action? Cicero's question remains: does foreknowledge cause the future, or does it only describe what will happen?

\subsection*{8. Cicero's Challenge}
Cicero does not offer a simple slogan. He sees that human actions have causes. He also sees that moral life collapses if causes become excuses for everything.

His challenge is to distinguish cause from compulsion. Nature influences us. Circumstances pressure us. Character shapes what seems good. But if we call all of this fate in a way that destroys choice, then praise, blame, law, teaching, and repentance lose their meaning.

Cicero therefore resists strict fate. He wants room for human beings to deliberate, assent, refuse, learn, and be responsible. The hard question is how much room remains once we admit that every action is surrounded by causes.

\section*{Reading Focus}
\begin{enumerate}[leftmargin=*]
\item What is the difference between saying an event has a cause and saying it happened by fate?
\item Why does Cicero think nature and chance can explain many events without invoking fate?
\item Explain the lazy argument. Why does Cicero reject it?
\item What does the rolling cylinder image suggest about outside causes and inner character?
\item Cicero asks whether prediction causes the future. How might this question matter for \textit{Oedipus Rex} or \textit{Macbeth}?
\end{enumerate}

\section*{Window Note and Summary}
Create a window note showing three forces in Cicero's argument: nature, chance, and choice. Include at least one drawing or symbol for each. At the bottom, write a 3--5 sentence summary explaining why Cicero worries that strict fate would weaken praise, blame, law, and moral responsibility.

\section*{Connection to the Unit Question}
Write one claim answering this question: \textbf{If our choices have causes, does that mean they are fated?} Use Cicero to support your answer.
\end{multicols}
\end{document}
